% Chapter 5: Discussion
\chapter{Discussion}
\label{ch:discussion}

This chapter interprets the results from Chapter~\ref{ch:results} and discusses the methodological implications, limitations, and future research directions of MutBench.
The multi-source information integration analysis (Section~\ref{sec:information_analysis}) reinforces the cross-pathogen findings: multi-feature integration consistently outperformed single-scoring approaches, complementing the statistical evidence from the benchmark evaluation.

%% ============================================================
\section{Methodological Validity of MutBench}

%% ============================================================
\subsection{Interpretation of Bootstrap Confidence Intervals}
\label{sec:bootstrap_interpretation}

The wide 95\% confidence intervals observed in region-level BCa bootstrap (e.g., MutClust-Hybrid [0.323, 0.635]) reflect region-dependent recall sensitivity in small samples ($n = 7$): resamples dominated by high-signal regions (RBD) inflate recall, while low-signal-only resamples deflate it.
Nevertheless, pairwise bootstrap comparisons show that differences between methods are significant ($p < 0.001$), because paired evaluation on the same resample cancels region-combination variability.
Future work should expand the effective sample size by including additional protein domains (e.g., ORF1ab, N protein).
Having established statistical robustness, the next critical question is whether synthetic benchmark conclusions transfer to real data.

%% ============================================================
\subsection{Performance Gap Between Synthetic and Real Data}
\label{sec:synthetic_real_gap}

The real GISAID benchmark (Section~\ref{sec:real_gisaid_benchmark}) demonstrated that the sparse distribution of real H-scores (342/29,903 nonzero) fundamentally changes detection dynamics: MutClust-Hybrid's synthetic advantage disappears (hotspot-score 0.383 vs.\ 0.386 for MutClust-Fixed).
However, enrichment ratios remain 6.9--8.0-fold above random, confirming that biological signal is preserved despite the F1 drop.
The dual-benchmark design---quantifying the gap between theoretical detection capacity (synthetic) and practical performance (real)---is itself a methodological contribution, and MutBench recommends evaluating both as complementary perspectives.

%% ============================================================
\subsection{Independence of Ground Truth and Detection Methods}
\label{sec:circularity}

The synthetic benchmark's ground truth is defined independently of the detection methods (DMS experimental data and convergent evolution literature), ensuring no \emph{algorithmic} circularity between evaluation and detection.
A more subtle concern is that Layer~A positions (convergent evolution) could correlate with mutation frequency at the \emph{per-position} level, which can make frequency-based scoring appear structurally strong. A per-pathogen audit of the point-biserial correlation between Layer~A membership and per-position frequency finds a mean $|\rho| = 0.12$ across the 11 pathogens (range $[-0.06, +0.33]$). This per-position membership-vs-frequency $|\rho|=0.12$ audit is conceptually distinct from the feature-feature Pearson correlation $\rho \approx 0.97$ between the \emph{frequency} and \emph{entropy} scoring inputs reported in Section~\ref{sec:information_analysis} (Chapter~\ref{ch:results}); the former measures circularity between ground truth and a scoring signal, while the latter measures redundancy among scoring inputs themselves. Per pathogen: HCV (+0.33), Influenza~B (+0.26) and Norovirus (+0.23) show moderate positive correlations, while SARS-CoV-2 ($\rho \approx 0$), MERS ($\approx 0$) and EV-A71 ($\approx 0$) are essentially uncorrelated. The effect is therefore pathogen-specific rather than uniform, and is driven by ground-truth curation style: pathogens whose Layer~A uses hypervariable/epochal-variant criteria (HCV, Influenza~B, Norovirus) couple more strongly with frequency than pathogens whose Layer~A uses phylogenetically-confirmed convergent positions (SARS-CoV-2, H3N2).
At the individual-position level, D614G achieved the highest single frequency through a founder event (homoplasy count = 2), while position~476 has moderate frequency but the highest homoplasy count (53); such anecdotes confirm that frequency and convergent-evolution signals are not interchangeable. Readers should therefore interpret frequency scoring's strong benchmark performance as partly structural (frequency $\approx$ Layer~A by construction) and partly genuine signal (homoplasy provides the 3-lineage discriminator that mere frequency cannot).
To provide a frequency-independent validation dimension, Layer~C (DMS experimental fitness) is used where available (6 of 11 pathogens), and the vaccine escape enrichment analysis (Section~\ref{sec:vaccine_escape}) provides partial external cross-validation: HIV-1 (82\% of its escape set disjoint from Layer~A) is the cleanest external anchor, whereas H3N2's 69\% Layer~A/escape overlap means that the H3N2 enrichment re-measures benchmark performance more than it provides independent replication. The circularity audit in Chapter~\ref{ch:results} Table~\ref{tab:vaccine_escape_stage3} reports novel-only enrichment values that quantify this distinction.
The biological validity of detected hotspots was further supported by 3D structural analysis (Section~\ref{sec:structure_analysis}), which showed 5.92-fold spatial contact enrichment, and by the consistent superiority of domain-specific methods (MutClust variants) over general-purpose alternatives (OPTICS, KDE) in Section~\ref{sec:external_comparison}.
Min-max normalization was used as the default for detection method inputs; z-score normalization was applied only for the EqualWeight composite (Section~\ref{sec:scoring_detection}). Rank transformations were not evaluated as they can introduce artifacts in sparse data distributions.
With benchmark validity established, the temporal limitations of the underlying H-score metric warrant examination.

%% ============================================================
\section{Limitations of H-score-Based Detection}

%% ============================================================
\subsection{Limitations of Fixed-Reference Scoring}
\label{sec:hscore_saturation}

The H-score saturation phenomenon observed in temporal validation (Section~\ref{sec:temporal}) reveals a fundamental limitation of reference-based metrics: in 2024--2025 data, all 1,259 Spike positions (post-truncation alignment) have non-zero H-scores with a standard deviation of only 0.098, resulting in loss of discriminatory power.
TC-freq achieved the highest enrichment ratio (3.43-fold) on convergent evolution ground truth, suggesting that temporal contrast scoring is a viable direction; the appropriate temporal resolution remains a future challenge.
The saturation phenomenon motivates evaluation against multiple independent ground truth definitions.

%% ============================================================
\subsection{Complementary Evaluation Perspectives}
\label{sec:complementary_gt}

In the multiple ground truth evaluation (Section~\ref{sec:multi_gt}), the Jaccard similarity between DMS-escape and convergent evolution positions was only 0.006, confirming that the two measures capture fundamentally different aspects of variation.
Using only a single ground truth risks overestimating methods that align with a specific biological dimension, and MutBench's multi-ground-truth approach mitigates this.
The result that the best method varies depending on the ground truth suggests that the ``best hotspot detection method'' is context-dependent: for immune evasion studies, methods best suited for the convergent evolution criterion are preferred, while for molecular function studies, methods best suited for the DMS criterion are preferred.
MutBench systematically guides such context-dependent method selection.
The ground-truth-dependent method rankings observed within SARS-CoV-2 foreshadow the pathogen-dependent rankings revealed in the cross-pathogen benchmark.

%% ============================================================
\subsection{Homoplasy-Based Scoring as an Alternative}
\label{sec:homoplasy_scoring}

The founder bias problem ($\rho = -0.876$ between H-score and independent origin count; Section~\ref{subsec:phylo_correction_poc}) motivates homoplasy-based scoring as an alternative. Fitch parsimony-based homoplasy counting quantifies independent mutational origins on the phylogenetic tree, decoupling frequency from selective pressure: the D614G founder mutation receives a homoplasy count of only 2 (normalized 0.038), while convergent evolution position 476 received the highest count (53).
\phantomsection\label{tab:homoplasy_validation}

Cross-pathogen validation revealed pathogen-dependent performance: HIV-1 showed significant discrimination (AUC = 0.706, $p = 0.0002$), H3N2 showed a weak trend (AUC = 0.561, $p = 0.147$), and HCV failed entirely (AUC = 0.390) due to parsimony saturation in its hypervariable regions.
This failure mode implies that homoplasy-based scoring is unreliable for pathogens with extreme sequence diversity where parsimony inference breaks down. For such pathogens, frequency-based or structural scoring remains more appropriate.
This mirrors the detection-level finding of the main benchmark: the scoring $\times$ pathogen interaction ($\omega^2 = 0.296$) demonstrates that the optimal scoring strategy depends on pathogen-specific characteristics.

%% ============================================================
\section{Cross-Pathogen Generalization}

%% ============================================================
\subsection{Implications of Cross-Pathogen Generalization}
\label{sec:cross_pathogen_implications}

The cross-pathogen benchmark demonstrated that no universally best-performing method exists (Chapter~\ref{ch:results}), with implications for real-time surveillance systems such as Nextstrain~\cite{nextstrain2018hadfield}: applying a previously optimized combination to a new pathogen may not generalize.
The underlying cause is that pathogen-specific genomic characteristics create fundamentally different signal landscapes---for example, H3N2's uniformly nonzero H-scores neutralize threshold-based seed detection, while MERS's sparse genome favors weak-signal detectors.

A notable finding is the asymmetry between scoring and detection contributions: the detection method family main effect ($\omega^2 = 0.013$) is the smallest variance component, approximately 23 times smaller than the scoring $\times$ pathogen interaction ($\omega^2 = 0.296$). This implies that \textit{the choice of input information} matters far more than \textit{the choice of detection algorithm}. Practically, this means that effort should be directed toward selecting the right information source for each pathogen rather than optimizing detection algorithm parameters.

These findings are formalizable within Rice's algorithm selection problem~\cite{rice1976algorithm}: the problem space $\mathcal{P}$ consists of pathogens characterized by genomic profiles, the algorithm space $\mathcal{A}$ comprises 780 scoring--detection combinations, and the dominant interaction effect demonstrates that the mapping $f: \mathcal{P} \to \mathcal{A}$ is non-trivial---consistent with the No Free Lunch theorem~\cite{wolpert1997nfl}---though unlike the theoretical NFL result, which applies to all possible problem distributions, the empirical finding here is grounded in a specific, biologically meaningful problem class (11 RNA viruses with documented functional sites) and therefore provides stronger practical evidence for pathogen-adaptive selection than the abstract theorem alone.
The ESM-2 cross-paradigm validation (Section~\ref{subsec:esm_validation}) reinforces this: ESM-LLR is the top-ranked feature for SARS-CoV-2 but substantially outperformed by homoplasy for Norovirus (AUC $= 0.763$ vs.\ $0.944$), confirming that no single scoring paradigm dominates.

\subsection{Comparison with Prior Work}
\label{sec:comparison_prior_work}

The pathogen-dependence observed in MutBench is consistent with related work, but MutBench addresses a different task.
EVEREST~\cite{gurev2025everest} and ProteinGym~\cite{notin2023proteingym} show that variant-effect-prediction performance varies across proteins and viral families, while Maher et al.~\cite{maher2022predicting} and Hie et al.~\cite{hie2021learning} show that different feature types capture different aspects of SARS-CoV-2 or viral escape.
Those studies focus mainly on per-variant fitness or escape prediction; MutBench instead evaluates per-region hotspot detection across a curated 11-RNA-virus panel.
Its contribution is therefore the quantified scoring~$\times$~pathogen interaction ($\omega^2 = 0.296$ at 20-type granularity; 0.103 at 6-category granularity) and the practical cross-check that HIV-1 escape enrichment (7.19$\times$, 82\% disjoint from Layer~A) provides the strongest external anchor, whereas H3N2 is mainly self-consistency evidence.

The demonstrated necessity of pathogen-adaptive selection raises the question of how to interpret the absolute performance levels achieved.

%% ============================================================
\subsection{Upper Bound of Detection Performance}
\label{sec:oracle_mcc_ceiling}

The mean oracle MCC across 11 pathogens is 0.341 (range: 0.196--0.660; Table~\ref{tab:stage2_best}), which appears low but is explained by three structural factors: position-level exact matching (off-by-one counts as false positive), extreme class imbalance (positive rates 0.7\%--15.9\%), and ground truth stringency (only literature-confirmed sites).
The oracle is unambiguously non-random (113 standard deviations above the random baseline mean of 0.001), and comparable to position-level metrics in other domains (e.g., TFBS prediction F1 = 0.3--0.5~\cite{weirauch2013evaluation}).

Region-overlap analysis (Table~\ref{tab:region_overlap_mcc}) quantifies the effect of evaluation stringency: under $\pm5$ position windows, mean oracle MCC increases from 0.341 to 0.472 (38\% improvement), with Influenza~B reaching 0.894. Notably, some pathogens (HCV, Norovirus) show \textit{decreased} MCC under window expansion because their oracle detections are already precise---window expansion dilutes precision by including non-functional neighboring positions.

\begin{table}[htbp]
\centering
\caption{Region-overlap MCC under three evaluation modes. Upper block (9 pathogens) reproduces directly from \texttt{region\_overlap\_mcc.csv}; lower block (EV-A71, Rabies) uses oracle-combination derivation (fixed-combination CSV rows pending regeneration). Window columns: $\pm 5$, $\pm 10$ position tolerance.}
\label{tab:region_overlap_mcc}
\small
\begin{tabular}{lccc}
\toprule
\textbf{Pathogen} & \textbf{Exact} & \textbf{$\pm$5 window} & \textbf{$\pm$10 window} \\
\midrule
\multicolumn{4}{l}{\textit{Reproducible from \texttt{region\_overlap\_mcc.csv} (9 pathogens):}} \\
HCV           & 0.660 & 0.588 & 0.519 \\
H3N2          & 0.534 & 0.639 & 0.548 \\
Norovirus     & 0.510 & 0.335 & 0.275 \\
Influenza B   & 0.392 & 0.894 & 0.924 \\
HIV-1         & 0.275 & 0.406 & 0.343 \\
Dengue        & 0.274 & 0.441 & 0.508 \\
SARS-CoV-2    & 0.211 & 0.480 & 0.497 \\
MERS          & 0.196 & 0.470 & 0.593 \\
RSV           & 0.196 & 0.394 & 0.318 \\
\midrule
\textit{Subset mean (9)} & 0.361 & 0.516 & 0.503 \\
\midrule
\multicolumn{4}{l}{\textit{Oracle-derived (no archived fixed-combination CSV row):}} \\
EV-A71        & 0.260 & 0.275 & 0.199 \\
Rabies        & 0.248 & 0.265 & 0.266 \\
\midrule
\textbf{Mean (all 11)} & \textbf{0.341} & \textbf{0.472} & \textbf{0.454} \\
\bottomrule
\end{tabular}
\end{table}

Despite these ceiling effects, the benchmark results have practical implications for public health surveillance.

\subsection{Recall Analysis: Which Positions Are Missed?}
\label{sec:recall_analysis}

The mean oracle recall across 11 pathogens is 0.552 (Table~\ref{tab:precision_at_k}), indicating that the best combination for each pathogen detects approximately half of the known convergent evolution positions. While this may appear insufficient for safety-critical applications, three observations provide context.

First, \textit{the missed positions are not randomly distributed}. Positions with low recall tend to be those where the convergent evolution signal is weakest---positions documented in only 3 lineages (the minimum threshold for Layer~A inclusion), or positions where the selective advantage is epistatic (dependent on other mutations) rather than individually strong. For SARS-CoV-2, the 8 undetected Layer~A positions (recall = 0.680, 17/25 detected) are predominantly NTD supersite positions where immune evasion operates through conformational changes rather than direct frequency elevation.

Second, \textit{multi-source integration can recover positions missed by a single-information view}. In the H3N2 integration case study, frequency-based top-5 selection recovers a subset of antigenic site~B positions but misses others that are better prioritized by the EqualWeight ensemble through the pLDDT and homoplasy components (Section~\ref{subsec:search_space_reduction}). This complementarity---where different information types detect different subsets of positions---is the mechanistic basis for the dominant scoring $\times$ pathogen interaction.

Third, \textit{region-level coverage is substantially higher than position-level recall}. Under $\pm$10 position windows (approximately one epitope width), the mean MCC increases from 0.341 to 0.454, and Influenza~B reaches 0.924. Positions scored as ``missed'' in exact evaluation are often detected within 1--2 residues of the true position---close enough to guide experimental prioritization, which is the practical goal.

The primary risk of missed positions lies in pathogens with low recall: EV-A71 (recall = 0.185) and Rabies (recall = 0.308). For these pathogens, the 261~AA and 491~AA post-truncation alignment lengths provide limited statistical power, and the small number of detected positions ($n = 10$ and $n = 39$, respectively) may not adequately represent the full landscape of adaptive sites. Users applying MutBench to such short-protein pathogens should interpret the detected positions as a prioritized subset rather than an exhaustive list, and complement computational detection with experimental validation (DMS or epitope mapping).

%% ============================================================
\section{Practical Implications and Fundamental Limitations}

Returning to the three problems identified in Chapter~\ref{ch:introduction}: the standardised evaluation framework is provided by the three-layer ground truth and hotspot-score metric (Contribution~1); pathogen-dependence is quantified by the interaction effect ($\omega^2 = 0.296$) and corroborated by LOPO 0/11 (Contribution~2); cross-validation against vaccine-escape positions establishes practical relevance, anchored by HIV-1 (Contribution~3).

%% ============================================================
\subsection{Vaccine Escape Position Validation}
\label{sec:vaccine_escape}

Vaccine escape validation (Table~\ref{tab:vaccine_escape_stage3}) showed enrichment values of 9.36$\times$ for H3N2, 7.19$\times$ for HIV-1, and 7.63$\times$ for SARS-CoV-2 (nominal Fisher's exact $p$-values). Under Bonferroni correction across 780 combinations, HIV-1 ($p_{\text{adj}} = 2.5 \times 10^{-16}$) and H3N2 ($p_{\text{adj}} = 2.6 \times 10^{-5}$) survive while SARS-CoV-2 does not. The circularity audit in Chapter~\ref{ch:results} further distinguishes HIV-1 (82\% of escape set disjoint from Layer~A; external validation) from H3N2 (69\% Layer~A overlap; self-consistency check).
The 10-feature EqualWeight integration (top-5\% threshold; \texttt{escape\_validation\_adapt.csv}) further achieves statistically significant H3N2 escape enrichment (4.012$\times$, $p = 0.0023$) in the integration-vs-single comparison reported in that file (note: single-scoring H3N2 FUBAR survives Bonferroni in Table~\ref{tab:vaccine_escape_stage3}; the "single-feature approaches do not" statement here applies specifically to the EqualWeight/FreqThresh/BestSingle top-5 integration block, not to the full 780-combination Stage 2 sweep). The primary practical value lies in \textit{reducing the experimental search space} rather than achieving high absolute prediction accuracy (see Section~\ref{subsec:search_space_reduction} for detailed discussion).

While LOPO cross-validation yields 0/11 generalization, MutBench provides practical value through three mechanisms: (1)~the family-level scoring recommendation table (Table~\ref{tab:scoring_recommendation} in this chapter) provides coarse-grained starting hypotheses based on virus family characteristics; (2)~the 10-feature EqualWeight ensemble (4-feature core retains $\approx$98\% of its mean MCC) achieves 4.012-fold H3N2 vaccine escape enrichment in the integration comparison without requiring pathogen-specific optimization; and (3)~once scoring features have been computed for a pathogen, the 20-scoring $\times$ 39-detector evaluation grid can be rerun in approximately 25~seconds on a single workstation, making exhaustive detector comparison operationally feasible.
The earlier Stage~1 single-scoring escape analysis is therefore retained only as historical development context; the dissertation-level evidence uses the full 20-scoring benchmark in Table~\ref{tab:vaccine_escape_stage3}.

%% ============================================================
\subsection{Cross-Validation via Vaccine Escape Enrichment}
\label{sec:cross_validation_escape}

A critical question for any benchmark is whether its rankings reflect genuine biological relevance or merely circular evaluation artifacts.
The vaccine escape enrichment analysis across three pathogens (2{,}340 evaluations) provides \emph{partial} cross-validation of the Layer~A benchmark results, with the strength of external validation differing substantially by pathogen.

\paragraph{Three-tier evidence rule.}
External value is judged by Bonferroni survival and Layer~A independence.
HIV-1 passes both criteria (18\% Layer~A overlap, $p_{\text{adj}} = 2.5 \times 10^{-16}$) and remains the primary external anchor; the direct novel-only test on 37 Layer~A-disjoint positions gives $7.16\text{--}8.24\times$ enrichment, with the conservative worst case still significant after 780-fold correction.
H3N2 survives Bonferroni correction but has 69\% Layer~A overlap, so it is treated as a self-consistency check; its novel-only enrichment is not significant ($p=0.132$).
SARS-CoV-2 does not survive correction and is exploratory.
Thus the benchmark is not algorithmically circular, but the practical-value claim rests mainly on HIV-1.

\textbf{Concurrent-benchmark task-orthogonality.}
ViroGym and EVEREST overlap with MutBench on some pathogens, but they optimize per-variant fitness Spearman whereas MutBench evaluates per-region binary hotspot calls.
A fair head-to-head would require task-aligned binarization of fitness predictions, so it remains future work rather than a missing leaderboard.
This is especially relevant for pathogens lacking DMS Layer~C, where Layer~A plus HIV-1-like external escape validation remains the current evidential bridge.

%% ============================================================
\subsection{Practical Considerations for Public Health Surveillance}
\label{sec:practical_surveillance}

Given SARS-CoV-2's generation time ({$\sim$}5 days) and sequence submission delay ({$\sim$}14 days), a minimum 4-week window is required for real-time TC-freq alerts. A fundamental methodological limitation affecting all frequency-based approaches is phylogenetic non-independence.

For public health practitioners, MutBench results can be operationalized through the following workflow when a novel pathogen or variant emerges: (1)~compute the 20 scoring types from the available MSA and auxiliary inputs (MAFFT for alignment, HyPhy for FUBAR, ESM-2 for PLM scores, and AlphaFold/ESMFold or experimental structures when structural features are used); (2)~if the pathogen belongs to a family already represented in MutBench (e.g., Coronaviridae, Orthomyxoviridae, Flaviviridae), use the family-level scoring recommendation from Table~\ref{tab:scoring_recommendation} as a starting hypothesis rather than a fixed rule; (3)~if the pathogen is novel or labelled sites are unavailable, apply the EqualWeight ensemble as a conservative cold-start baseline, noting that the 4-feature core (homoplasy, pLDDT, entropy, frequency) retained 97.6\% of the full 10-feature ensemble's mean MCC under the specific uniform-weight LOPO top-10\% protocol. This workflow integrates with existing surveillance platforms such as Nextstrain~\cite{nextstrain2018hadfield} and GISAID~\cite{shu2017gisaid}, which already provide the MSA inputs required by MutBench.

\paragraph{Optimality ceiling vs robust default floor.}
Contribution~2 gives the informed-use ceiling: pathogen-specific selection explains the scoring~$\times$~pathogen interaction and the 0.265 oracle-vs-generalized gap.
Contribution~3a gives the cold-start floor: the 4-feature core recovers $\approx$98\% of the 10-feature ensemble mean MCC when labels are unavailable.
These are complementary operating regimes, not a single validated universal adaptive algorithm.

\paragraph{PAHD-R as a bounded synthesis, not a universal selector.}
PAHD-R narrows the gap between the ceiling and floor without erasing the limitation.
It should be interpreted as a shared evidence-fusion and review framework with three global modes: Core for no-AI exact/top-$k$ reporting, Augmented for balanced region prioritization, and Review for conservative low-FPR triage.
Calibrated, shrinkage, virus-aware, and callable-only variants are feasibility signals for adaptive extension, but they remain candidates because none simultaneously clears cross-pathogen robustness, per-pathogen null support, non-abstaining coverage, and independent validation.
Repair layers are also bounded: HCV E2 is adopted after coordinate reconciliation, while SARS-CoV-2 RBD/ACE2 and MERS DPP4/contact remain candidates and RSV remains unrepaired.

%% ============================================================
\subsection{Practical Application Workflow for New Pathogens}
\label{sec:practical_workflow}

Based on the cross-pathogen findings, we outline a concrete workflow for applying MutBench's methodology to a novel RNA virus not included in the benchmark.

\paragraph{Step 1: Sequence collection.}
Retrieve $\geq$200 protein-coding sequences from NCBI GenBank or GISAID, ensuring temporal and geographic diversity. Construct a multiple sequence alignment using MAFFT~\cite{katoh2013mafft} (${\sim}$5 minutes for 1{,}000 sequences).

\paragraph{Step 2: Feature computation.}
Compute scoring features in order of increasing computational cost: (i)~frequency and Shannon entropy from the MSA (seconds); (ii)~homoplasy counts via Fitch parsimony on an inferred phylogenetic tree (2--5 minutes in the benchmark-scale examples); (iii)~ESM-2 masked-marginal scores, requiring a GPU for practical throughput (3--5 minutes per protein in the benchmark environment); (iv)~structural features such as pLDDT and surface accessibility, requiring an AlphaFold/ESMFold predicted structure or an experimental PDB entry when available. Steps (i)--(iii) are broadly applicable to pathogens with sufficient sequence diversity; step (iv) can be omitted when no reliable structural model exists, with the resulting analysis interpreted as a reduced-feature baseline. When structural features are available, the 4-feature core (homoplasy, pLDDT, entropy, frequency) captures 98\% of full-ensemble performance; without pLDDT, the remaining 3 features (homoplasy, entropy, frequency) still provide the dominant scoring signals.

\paragraph{Step 3: Scoring category selection.}
Table~\ref{tab:scoring_recommendation} provides empirically grounded starting categories based on which scoring category achieved the highest mean MCC for each virus family in the benchmark. Users should treat the recommended category as a first-pass prior, not as a general algorithmic rule, and should evaluate additional categories whenever computational resources permit.

\begin{table}[htbp]
\centering
\caption{Starting scoring-category priors by virus family, based on cross-pathogen benchmark results}
\label{tab:scoring_recommendation}
\small
\begin{tabularx}{\textwidth}{llX}
\toprule
\textbf{Virus Family} & \textbf{Starting Scoring Prior} & \textbf{Rationale} \\
\midrule
Orthomyxoviridae & Phylogenetic (FUBAR) & Strong diversifying selection signal \\
Coronaviridae    & AI-based (PLM)       & Shallow phylogeny; PLM captures broader context \\
Caliciviridae    & MSA-derived (homoplasy) & Epochal evolution with convergent patterns \\
Retroviridae     & MSA-derived (entropy)   & High diversity; entropy captures allelic variation \\
Flaviviridae     & Frequency-based         & Hypervariable regions dominate signal \\
Pneumoviridae    & AI-based (ESM-2)        & Limited phylogenetic signal \\
\bottomrule
\end{tabularx}
\end{table}

\paragraph{Step 4: Detection and evaluation.}
Apply KDE or Wavelet-based detection as robust defaults (Section~\ref{sec:external_comparison}). If ground truth positions are available (e.g., from DMS experiments or literature), compute the hotspot-score metric. When scoring features are pre-computed, detection and evaluation across all 39 detector configurations completes in under 3~seconds per pathogen (the full 8,580-evaluation benchmark requires approximately 25~seconds).

\paragraph{Computational requirements.}
The total pipeline runtime per pathogen is approximately 15--20 minutes on a workstation equipped with a single GPU (tested on NVIDIA RTX 3090). The most time-consuming step is ESM-2 inference; all other steps complete in under 10 minutes on a standard CPU. No specialized high-performance computing infrastructure is required.
For resource-constrained settings, the 4-feature core (homoplasy, pLDDT, entropy, frequency) numerically retains 98\% of the full 10-feature ensemble mean MCC on the 12-pathogen Stage~3 panel under uniform-weight LOPO top-10\% evaluation (0.081 vs.\ 0.083; no paired-sample significance test was performed, so the gap is reported as a numerical rather than statistical equivalence claim), while requiring only MSA-based computation plus a single AlphaFold2 structure query---eliminating PLM inference (no ESM-2 forward passes) and bringing runtime to under 10 minutes per pathogen on a CPU. Note that AlphaFold2 itself remains compute-intensive (typically GPU-accelerated and on the order of tens of minutes per protein in reference setups); the 4-feature core saves on PLM inference, not on all heavy compute, and would benefit further if cached AlphaFold-DB structures are reused instead of running fresh predictions.

%% ============================================================
\subsection{Phylogenetic Non-Independence and Scalability Limits}
\label{sec:phylo_limits}

All detection methods in this study implicitly assume independence between sequences; however, SARS-CoV-2 sequences are phylogenetically structured, and frequency-based approaches cannot distinguish positive selection from founder effects (e.g., D614G;~\cite{korber2020tracking}).
The dN/dS filter partially addresses this but is itself susceptible to phylogenetic bias.

Proof-of-concept simulations (Section~\ref{subsec:phylo_correction_poc}) demonstrated that Fitch parsimony-based homoplasy scores correctly rank convergent mutations above founder mutations, and extended TreeTime ML validation (Section~\ref{subsec:treetime_ml_validation}) confirmed equivalent accuracy (AUC = 1.000).
Validation on real GISAID sequences (Section~\ref{subsec:treetime_gisaid_validation}) showed that H-score and homoplasy score capture fundamentally different signals ($\rho = -0.107$, $p = 1.36 \times 10^{-4}$), supporting the practical value of homoplasy-based correction.
The homoplasy sensitivity analysis (Section~\ref{subsec:homoplasy_sensitivity}) revealed a temporal mismatch between the 2020-H1 sequence data and the Layer~A convergent evolution ground truth, confirming that phylogenetic correction requires temporally aligned data.
Full integration of phylogenetic correction with temporally appropriate sequence sampling remains a key future extension (Chapter~\ref{ch:conclusion}).

%% ============================================================
\subsection{Representativeness and Sampling Bias of GISAID Data}

Systematic biases exist in the GISAID~\cite{shu2017gisaid}-based input data.
\textbf{Geographic bias}: High-income countries contribute over 80\% of total sequences~\cite{chen2022gisaid}, leading to underrepresentation of unique variants from Africa/South Asia; weekly stratification corrects only for temporal imbalance.
\textbf{Temporal bias}: Sequence accumulation differs substantially between the early pandemic and Omicron periods, and within-week representativeness is not guaranteed.
\textbf{Founder effect}: ``Hitchhiker'' mutations accompanying lineage expansion (e.g., JN.1;~\cite{kaku2024virological}) cannot be distinguished from positively selected positions by frequency-based approaches alone.

\textbf{Temporal window dependence}: An additional unstandardized factor is the analysis time window. The 11 pathogens in this benchmark were analyzed using sequences spanning different time ranges: SARS-CoV-2 covers approximately 2 years of pandemic evolution (2020--2021), while H3N2 encompasses decades of seasonal circulation.
Temporal analysis on SARS-CoV-2 Spike across three 6-month windows reveals that feature discriminative power changes substantially: frequency AUC increases from 0.337 (2020-H1) to 0.472 (2021-H1), and entropy AUC from 0.328 to 0.507, as mean Shannon entropy rises 10-fold (0.126 to 1.241) with accumulating diversity.
This implies that the optimal scoring type for a given pathogen may depend not only on its biology but also on the temporal depth of the available sequence data. Standardizing the analysis time window---or explicitly studying how optimal scoring changes with temporal coverage---is an important future direction for hotspot benchmarking.

\subsection{ESM-2 Training Data Overlap}
\label{sec:esm2_leakage}

ESM-2~\cite{lin2023esm2} was trained on UniRef50/UniRef90, which contains viral protein sequences, and the 11 target proteins in our benchmark (e.g., SARS-CoV-2 Spike, HIV-1 gp120, H3N2 HA) are likely present in the training corpus.
Consequently, ESM-2 masked-marginal scores may partially reflect memorized evolutionary patterns rather than purely generalizable protein language understanding, potentially inflating absolute AUC values.
However, this concern applies uniformly to all 11 pathogens in the analysis, so while absolute performance levels may be biased upward, the relative cross-pathogen comparisons---which constitute our main finding---remain valid.
The semantic change feature (cosine distance between wild-type and mutant embeddings) is less susceptible to this concern because it measures relative embedding shifts rather than absolute likelihoods, making it more robust to training set memorization.
Future work should evaluate protein language models trained on databases that exclude viral sequences to quantify the magnitude of this potential bias. \textbf{Concurrent-benchmark reconciliation}: EVEREST v3 reports that PLMs underperform alignment-based methods on viral DMS \emph{fitness regression}, while the present work finds PLMs winning on 2 of 11 pathogens for region-level \emph{detection} (RSV best ESM-2 masked-marginal channel; SARS-CoV-2 best Tranception). The two findings are compatible because the tasks differ structurally (per-variant continuous Spearman vs.\ per-region binary MCC), and the fact that Norovirus's best is homoplasy (AUC 0.944) rather than a PLM is exactly the pattern leakage-driven inflation would not predict.

\textbf{Detector-family Simpson's-paradox check}: the family main effect ($\omega^2 = 0.013$) is the smallest modeled component, but the family~$\times$~pathogen interaction (0.082, medium) shows that family choice does matter when conditioned on pathogen. The within-cell residual (0.39) absorbs the unmodeled three-way interaction and parameter-variant noise; a fully-resolved four-way ANOVA with parameter-variant as a fourth factor would localise that residual but is deferred (computational cost; interpretability of $20 \times 14 \times 11 \times 39 \sim 10^5$ cell means is poor). The framing ``information matters more than algorithm'' compares modeled main effects only and is not a claim that detection algorithms are irrelevant within a fixed (scoring, pathogen) cell.

Among the 20 scoring types, several share an underlying information source (e.g., freq, freq\_with\_indel, and rare\_freq are all frequency-derived). This partial collinearity could inflate the apparent number of independent scoring levels in the ANOVA. When scoring types are aggregated to the 6-category level, the scoring $\times$ pathogen interaction shrinks substantially: a supplementary two-way ANOVA with 6 scoring categories yields $\omega^2_{\text{category} \times \text{pathogen}} = 0.103$, which is below Cohen's large-effect threshold ($\omega^2 \geq 0.14$) but above the medium threshold ($\omega^2 \geq 0.06$), and is comparable in magnitude to the pathogen main effect ($\omega^2 = 0.117$). The interaction therefore remains practically non-trivial even after collinearity correction, but the ``largest source of variance'' framing in the 20-type model is partly driven by fine-grained distinctions within categories; readers should interpret the 6-category $\omega^2 = 0.103$ as the collinearity-robust lower bound and the 20-type $\omega^2 = 0.296$ as the upper bound on the pathogen-scoring interaction depending on how ``scoring type'' is defined.

A methodological limitation concerns the \emph{winner's curse}~\cite{ioannidis2008winner}: with 780 candidate combinations searched per pathogen and the highest-MCC combination retained, the reported top-1 enrichment is an upper-envelope estimate rather than a typical value (the Bonferroni-survivor mean enrichment is materially lower; H3N2 mean $\approx 4.7\times$ vs.\ top $9.36\times$). The winner's-curse formulation is one face of the broader question of multiple testing correction when selecting the best scoring--detection combination per pathogen. At the detection-family level (20 scoring $\times$ 14 families = 280 combinations evaluated per pathogen), the probability that 11 out of 11 pathogens exhibit unique best combinations by chance alone is approximately 82\% under uniform random assignment ($\prod_{i=0}^{10} (280-i)/280 \approx 0.820$); at the parameter-variant level (20 scoring $\times$ 39 variants = 780 combinations, as reported in Section~\ref{subsec:9pathogen_statistical}), the same probability is approximately 93\%. The choice of aggregation level matters: the family-level calculation is the more conservative reference for cross-pathogen uniqueness because parameter-variant collapsing to families reduces the effective combinatorial space. However, several factors mitigate this concern: (1)~the best combinations are not merely unique but involve 9~distinct scoring \textit{types} from 4~of~6 categories, which has much lower random probability; (2)~the ANOVA interaction effect ($\omega^2 = 0.296$) quantifies systematic rather than chance variation; and (3)~the vaccine escape enrichment cross-validation confirms that benchmark-optimal methods also achieve the highest enrichment on independently curated escape positions. Nevertheless, future work should apply permutation-based significance testing to determine whether each pathogen's best combination is statistically distinguishable from alternatives.

The three-way ANOVA treats pathogen as a fixed effect, which is appropriate for inference about these 11~specific pathogens but does not support generalization to all RNA viruses. A mixed-effects model treating pathogen as a random effect would enable broader inference; however, with only 11~pathogen-level observations, frequentist random-effect estimation is unreliable (Bolker et al.~\cite{bolker2009glmm} recommend $\geq$5--6 levels per random factor, ideally more). A natural alternative for the small-$n$ regime is \emph{Bayesian hierarchical (partial-pooling) modeling}, which we have additionally fit. The original exploratory calculation, with half-Cauchy priors on the between-scoring, between-pathogen, between-family, interaction-, and residual-scale parameters and a vague Normal prior on the grand mean (PyMC 5.28; 2~chains $\times$ 800 draws after 800 tuning steps; $\hat R \leq 1.01$), yielded a posterior mean of $\omega^2_{\text{scoring} \times \text{pathogen}} = 0.252$ with 95\% HDI $[0.188, 0.314]$ and posterior probability $\Pr(\omega^2_{\text{int}} > 0.14) = 0.996$. An archived re-implementation (\texttt{scripts/bayesian\_omega.py}; half-Normal $\sigma_{\bullet}$ priors with scale 0.2, evaluation-level model on the 8{,}580 observations, 4~chains $\times$ 2{,}000 draws after 2{,}000 tuning steps; \texttt{random\_seed=42}) gives a posterior mean of 0.319 with 95\% HDI $[0.246, 0.396]$ and $\Pr(\omega^2_{\text{int}} > 0.14) \approx 0.9998$. Both runs concur on the load-bearing posterior probability ($\Pr > 0.14$ near 1) and on the qualitative conclusion (large interaction); the difference in HDI center reflects the prior parameterization and aggregation level, and the archived run is provided as a reproducible provenance record (\texttt{results/mutbench/bayesian\_omega\_summary.csv}, \texttt{bayesian\_omega\_posterior.csv}). The partial-pooling estimate sits slightly below the fixed-effect point (0.296) as expected from variance shrinkage, but the credible-interval lower bound (0.188) still essentially coincides with the cluster-bootstrap lower bound (0.195) and remains above Cohen's large-effect threshold ($\omega^2 > 0.14$) under both inferential frames. As the benchmark expands to 20--30+ pathogens, frequentist mixed-effects modeling additionally becomes feasible by the Bolker criterion.

\textbf{Independent reproduction.} The 8{,}580-evaluation benchmark was produced by a single team using a single codebase against a single GISAID snapshot, with random seed~$=42$ and fully scripted configuration (Section~\ref{sec:code_availability}). Independent multi-tool reproduction on the same pathogen panel is queued for the next cycle (Priority~5)---a gap analogous to that identified by Bailey et al.~\cite{bailey2018comprehensive} prior to the PanCancer reproducibility effort.

\textbf{Technical-covariate sensitivity.}
Per-pathogen oracle MCC correlates strongly with two technical covariates that are not modeled in the three-way ANOVA: protein length (Spearman $\rho = -0.73$, $p = 0.011$ across the 11 pathogens) and Layer~A positive rate (Spearman $\rho = +0.64$, $p = 0.035$). That is, shorter proteins with a higher Layer~A positive rate (e.g., H3N2 HA, Norovirus capsid) systematically attain higher oracle MCC than longer proteins with sparse Layer~A (e.g., HCV E1E2). Because MCC itself is bounded by class balance, a portion of the reported $\omega^2_{\text{scoring} \times \text{pathogen}} = 0.296$ may reflect this structural heterogeneity rather than intrinsic biological signal. A formal ANCOVA fit (statsmodels Type~II OLS on the cell-level grid with $\log(\text{length})$ and Layer~A positive rate as continuous covariates) confirms the directional prediction: the pathogen main effect collapses from $\omega^2 = 0.108$ to essentially zero (covariates absorb it entirely, as expected since both are pathogen-constant), and the scoring~$\times$~pathogen interaction \emph{rises} from $\omega^2 = 0.234$ at the cell-level baseline to $\omega^2 = 0.264$ after adjustment ($p < 10^{-167}$ in both fits). Family~$\times$~pathogen also rises modestly (0.081 $\to$ 0.091), and scoring/family main effects are unchanged. The fixed-effect headline of 0.296 reported at the evaluation level (8{,}580 observations) is mechanically slightly higher than the cell-level 0.234 because aggregating to family-level mean MCC reduces within-cell residual; the \emph{relative shift} from baseline to ANCOVA-adjusted ($+0.030$) is the load-bearing quantity, and it confirms that the interaction is amplified rather than attenuated by technical-covariate adjustment.

\textbf{Layer-A curation-protocol subset analysis.} A second concern raised by the curation heterogeneity (Section~\ref{sec:gt_heterogeneity}) is that the interaction might be inflated by the most idiosyncratic Layer~A definitions (HCV's HVR diversifying selection and Norovirus's epochal-variant positions). Re-fitting the cell-level ANOVA on the HCV-excluded subset ($n = 10$) yields $\omega^2_{\text{scoring} \times \text{pathogen}} = 0.199$; further excluding Norovirus ($n = 9$) yields $\omega^2 = 0.187$. Both restricted estimates remain firmly in Cohen's medium-to-large range, demonstrating that the pathogen-adaptive finding is not driven by the two pathogens with non-convergent-evolution Layer~A construction. Restricting further to the four phylogenetically-derived Layer~A pathogens (SARS-CoV-2, H3N2, MERS, EV-A71; $n = 4$) gives $\omega^2 = 0.137$, still above Cohen's medium threshold despite the loss of statistical power. The pathogen-adaptive conclusion should therefore be read as holding \emph{within the 20-type/39-variant design and the 11-pathogen sample}, and independent validation on additional pathogens with matched length/positive-rate is a priority for follow-up work.

\phantomsection\label{sec:scoring_input_heterogeneity}\textbf{Scoring input heterogeneity.}
The 20 scoring types do not operate under identical conditions. Frequency and entropy scores require only the MSA, making them the most standardized. Phylogenetic scores (FUBAR, dN/dS, homoplasy) additionally depend on phylogenetic tree quality, which varies across pathogens: H3N2's well-sampled multi-decade tree provides more reliable dN/dS estimates than MERS's sparse 5-year phylogeny. Structure-based scores (pLDDT, SASA) depend on predicted structure accuracy, which is higher for well-characterized proteins (SARS-CoV-2 Spike) than for less-studied ones (EV-A71 VP1). AI-based scores inherit biases from their training data, as discussed above.
This input heterogeneity means that a scoring type's performance reflects not only the informativeness of its underlying biological signal but also the quality of auxiliary inputs (trees, structures, training data) available for each pathogen. The scoring $\times$ pathogen interaction therefore captures both genuine biological signal differences and input quality differences, which cannot be fully separated without controlled experiments using standardized auxiliary inputs.

However, a mitigating factor is that most detection methods used in this benchmark employ internal scale normalization. FreqThresh and AdaptiveKnee use percentile-based thresholds, KDE applies bandwidth-adaptive density estimation with percentile cutoffs, SlidingTest uses $t$-statistics, and SWAN applies within-window $z$-normalization. These mechanisms render the detection step largely scale-invariant, meaning that the raw scale differences across scoring types (e.g., frequency in $[0,1]$ versus homoplasy counts in $[0,53]$) do not systematically bias the comparison. The remaining scale-sensitive detectors (Bayes, ScoreDBSCAN) constitute a minority of the 14 families.

\subsection{Ground Truth Heterogeneity Across Pathogens}
\label{sec:gt_heterogeneity}

A key limitation of the 3-layer ground truth framework is the heterogeneity in how Layer~A (positive selection sites) is defined across pathogens.
For SARS-CoV-2 and H3N2, Layer~A positions are phylogenetically confirmed convergent evolution sites---positions where identical mutations arose independently in multiple lineages.
In contrast, HCV Layer~A comprises hypervariable regions (HVR1/HVR2) that exhibit diversifying rather than convergent selection, Norovirus Layer~A consists of epochal variant positions associated with antigenic shift events, and the remaining pathogens use literature-reported immune evasion sites with varying levels of phylogenetic confirmation (Table~\ref{tab:gt_positions}).
This heterogeneity could artificially inflate the method $\times$ pathogen interaction in the ANOVA, because differences in ground truth definition criteria---not only differences in pathogen biology---may contribute to the observed interaction effect.
However, this heterogeneity reflects the real-world challenge of hotspot detection: ground truth definitions necessarily depend on pathogen-specific biology, and no single criterion (e.g., convergent evolution alone) is applicable across all pathogens.
The finding that detection performance is pathogen-dependent holds regardless of whether this dependence stems from intrinsic biological differences among pathogens, from differences in ground truth construction, or from their interaction---in all cases, the practical implication is the same: a single detection method cannot be applied uniformly across pathogens without performance degradation.

The heterogeneous Layer~A definitions across pathogens (convergent evolution for SARS-CoV-2/H3N2, epochal variants for Norovirus, hypervariable regions for HCV) mean that MCC values are not directly comparable across pathogens in absolute terms. The ANOVA interaction effect should therefore be interpreted as reflecting systematic variation in \textit{relative} method rankings across pathogens, rather than absolute performance differences. The consistency of this finding under category-level aggregation ($\omega^2_{\text{category} \times \text{pathogen}} = 0.103$) provides additional robustness.

Additionally, HCV has an effective Layer~B size of 0 (constrained regions could not be mapped within the alignment length), meaning HCV evaluation relies solely on Layer~A. This asymmetry should be considered when interpreting HCV-specific results.
Future work could address this by incorporating structurally constrained positions from experimental HCV E2 structures or conservation analysis across Hepacivirus species.

\subsection{Scope Justification: 11 RNA Viruses}

The 11 pathogens were selected to systematically cover 8 major RNA virus families (Coronaviridae, Orthomyxoviridae, Retroviridae, Pneumoviridae, Flaviviridae, Caliciviridae, Rhabdoviridae, Picornaviridae), representing the primary taxonomic diversity of RNA viruses with pandemic potential.
This scope limitation to RNA viruses is an intentional design decision: H-score-based scoring requires sufficient mutation density, which RNA viruses provide due to their high mutation rates ($10^{-3}$--$10^{-5}$ substitutions/site/year), whereas DNA viruses ($10^{-6}$--$10^{-8}$) would require fundamentally different scoring approaches.
Despite the small sample size ($n=11$), the Friedman test yielded $p=0.990$, indicating that no single method can be reliably distinguished as superior across all pathogens. This is consistent with the LOPO finding (0/11 match rate), but because the LOPO outcome is null-consistent under random-combination permutation, the stronger evidence for pathogen-dependent method ranking remains the ANOVA interaction and the oracle-vs-generalized MCC gap. \textbf{Sampling-quality vs.\ pathogen-dependence}: the smallest-sample pathogens (Rabies $n_{\text{seq}} = 2{,}038$ but truncated 491 AA; EV-A71 $n_{\text{seq}} = 662$, 261 AA) do attain the lowest oracle MCCs (0.248, 0.260), so a sampling-quality confound is not zero. However, the technical-covariate sensitivity analysis (Section~\ref{sec:scoring_input_heterogeneity}, $\rho_L = -0.73$, $\rho_P = +0.64$) shows the joint covariate effect is subsumed by the pathogen main effect ($\omega^2 = 0.117$), and the scoring~$\times$~pathogen interaction (0.296) is a within-pathogen quantity that survives this absorption (see the ANCOVA bound there). Excluding the two smallest-sample pathogens (re-running the ANOVA on $n = 9$) shifts the headline interaction by less than the cluster-bootstrap CI half-width, so the pathogen-dependence finding is not a small-sample artifact of the two weakest panels.

The LOPO result should be interpreted with caution given the small sample size ($n = 11$). With only 10 training pathogens per fold, the power to detect a generalizable best combination is inherently limited. The 0/11 match rate is consistent with two interpretations: (1)~truly pathogen-specific optimality with no cross-pathogen generalization, or (2)~insufficient sample size to identify a robust generalizable pattern. The large mean performance gap (0.265 MCC) between oracle and generalized predictions supports interpretation~(1), but a definitive distinction requires a larger pathogen panel.

A methodological caveat is that the 11 pathogens are not phylogenetically independent: three pairs share viral families (SARS-CoV-2 and MERS in Coronaviridae, H3N2 and Influenza~B in Orthomyxoviridae, HCV and Dengue in Flaviviridae). The ANOVA treats pathogens as independent blocks, which may underestimate standard errors for pathogen-related effects. However, within each family pair, the best-performing scoring types differ (FUBAR Bayes factor for H3N2 vs.\ frequency for Influenza~B; frequency for HCV vs.\ entropy\_with\_indel for Dengue; Tranception for SARS-CoV-2 vs.\ frequency for MERS), indicating that phylogenetic relatedness does not determine method preference. This within-family heterogeneity strengthens rather than weakens the pathogen-adaptive finding.

\subsection{PLM and Data Limitations}
\label{sec:plm_data_limitations}

The AI-based scoring types (ESM-2 masked-marginal channel, Tranception proxy, semantic change, stability predictors) were evaluated using pre-trained models without virus-specific fine-tuning. The performance of these models is constrained by the composition of their training data: ESM-2 was trained on UniRef50 (general proteins), which includes limited representation of rapidly evolving viral surface proteins. Fine-tuning on viral protein families or using virus-specific language models could potentially improve performance but was not evaluated in this study. Recent work on CoVFit~\cite{covfit2025}, which fine-tunes ESM-2 on SARS-CoV-2 epidemiological and DMS data, demonstrates that pathogen-specific adaptation of PLMs can substantially improve fitness prediction---consistent with MutBench's central finding that the optimal information source is pathogen-dependent. Additionally, the comparison between sequence-based (ESM-2) and structure-conditioned (stability\_structural) predictors does not control for the quality of input structures, which varies across pathogens depending on the availability of experimentally determined structures versus AlphaFold/ESMFold predictions.
This may introduce systematic bias: pathogens with high-quality experimental structures (e.g., SARS-CoV-2 Spike) may show inflated structural feature importance relative to pathogens relying on predicted structures. Controlling for structure source (experimental vs.\ predicted) in future analyses would strengthen the structural feature comparisons.

The MSA sequences used in this benchmark were downloaded from NCBI GenBank without controlling for geographic or temporal sampling bias. Surveillance intensity varies dramatically across pathogens and regions: SARS-CoV-2 sequences are dominated by high-income country submissions during 2020--2022, while Rabies and EV-A71 sequences reflect more geographically diverse but temporally sparse sampling. This sampling heterogeneity may inflate mutation frequencies at founder-effect positions in over-sampled lineages (e.g., SARS-CoV-2 D614G) while underrepresenting diversity in under-sampled regions. The Layer~A ground truth partially mitigates this concern by requiring independent emergence across multiple lineages, but Layer~A position counts themselves may be biased toward well-studied pathogens.

\bigskip
\noindent\textbf{Chapter summary.}
This chapter discussed the implications of the MutBench results, including the pathogen-dependent nature of optimal information sources, the complementary role of multi-source integration, the limitations of the current framework, and the scope justification for 11 RNA viruses.
The following chapter summarizes the research contributions, analyzes remaining limitations, and presents future research directions.
